%%%%%%%% GenBio @ ICML 2026 — Workshop Draft v3 %%%%%%%%%%%%%%%%%
% v3: Complex Signals framing, softened claims, explicit hV4 logic

\documentclass{article}

\usepackage{microtype}
\usepackage{graphicx}
\usepackage{subcaption}
\usepackage{booktabs}
\usepackage{hyperref}

\newcommand{\theHalgorithm}{\arabic{algorithm}}

\usepackage{icml2026}

\usepackage{amsmath}
\usepackage{amssymb}
\usepackage{mathtools}
\usepackage{xcolor}
\usepackage{multirow}

\usepackage[capitalize,noabbrev]{cleveref}

% Placeholder command for figures not yet created
\newcommand{\figplaceholder}[2]{%
  \begin{center}
    \fbox{\parbox{#1}{\centering\vspace{1.5em}\textcolor{gray}{\textbf{[FIGURE PLACEHOLDER]}\\#2}\vspace{1.5em}}}
  \end{center}
}

\icmltitlerunning{Personalized Distortion Modeling in High-Dimensional Neural Color Representations}

\begin{document}

\twocolumn[
  \icmltitle{Inferring Subject-Specific Distortion Fields from \\
    High-Dimensional Neural Color Representations}

  \icmlsetsymbol{equal}{*}

  \begin{icmlauthorlist}
    \icmlauthor{Anonymous Authors}{anon}
  \end{icmlauthorlist}

  \icmlaffiliation{anon}{Anonymous Institution}

  \icmlcorrespondingauthor{Anonymous Author}{anonymous@example.com}

  \icmlkeywords{color vision deficiency, forward encoding, fMRI, personalized modeling, representational geometry}

  \vskip 0.3in
]

\printAffiliationsAndNotice{}

%==============================================================================
\begin{abstract}
%==============================================================================
Color vision deficiency (CVD) distorts the continuous geometry of cortical color representations, but characterizing this distortion from high-dimensional fMRI signals and mapping it to low-dimensional mechanistic parameters remains challenging.
We introduce a framework for fitting parsimonious, retina-inspired distortion models (1--2 free parameters) to individual fMRI-measured hue interpolation profiles in visual cortex.
Applied to two CVD and one control participant, the framework reveals that each CVD case is best captured by a \emph{different} model---a 1-parameter retinal cone-shift for moderate protanomaly and a 2-parameter retinal--cortical model for mild deuteranomaly---while the control is correctly rejected by all models.
The fitted models further predict whether stimulus-space correction is feasible: the mild case admits exact inversion, whereas the moderate case exhibits irreversible geometric collapse under the best-fitting model.
These results demonstrate that subject-specific distortion fields can be inferred from complex neural signals, with the degree of geometric collapse determining correction limits.
\end{abstract}

%==============================================================================
\section{Introduction}
\label{sec:intro}
%==============================================================================

Color vision deficiency (CVD) arises from altered spectral sensitivity of retinal cone photoreceptors and affects approximately 8\% of males worldwide \cite{placeholder_machado2009}.
Current assistive approaches, including commercial spectral-notch filters, apply a fixed, generic correction.
Recent evidence indicates that such filters can alter color \emph{appearance} but have minimal effect on color \emph{discrimination} at threshold \cite{placeholder_somers2024}, suggesting that effective correction may require individually-tailored strategies informed by each person's neural color geometry.

A central challenge is that CVD-related neural distortion is not uniform across individuals.
Even within a single diagnostic category (e.g., deuteranomaly), the degree of cortical compensation varies substantially \cite{placeholder_tregillus2021}, and the functional deficit that matters for continuous color tasks---the ability to interpolate between adjacent hues, not merely classify them---differs in ways that categorical diagnosis cannot capture.
Characterizing this distortion requires moving from classification-based neural decoding to \emph{geometry-sensitive} evaluation of high-dimensional fMRI representations: asking not whether the brain can tell two colors apart, but whether the continuous structure between them is preserved.

Here we propose a framework for inferring subject-specific distortion fields from fMRI data. More generally, we cast CVD as a case study in recovering low-dimensional distortion structure from complex representational geometry in high-dimensional neural data. We fit parsimonious mechanistic forward models to each CVD participant's neural interpolation profile, as measured by a geometry-sensitive leave-one-color-out (LOCO) cross-validation metric.
Our contributions are:
\begin{enumerate}
    \item We identify subject-specific interpolation failure as a geometry-sensitive phenotype that is more informative than coarse categorical decoding for characterizing CVD in our data, with hV4 providing the strongest signal for this phenotype.
    \item We introduce a personalized distortion-fitting framework in which different CVD participants are best captured by different parsimonious models (1--2 DOF), with successful specificity rejection of a normal control.
    \item We demonstrate that correction feasibility depends on the degree of geometric collapse induced by the fitted forward model, revealing severity-dependent limits on exact stimulus-space intervention.
\end{enumerate}

%==============================================================================
\section{Methods}
\label{sec:methods}
%==============================================================================

\subsection{Participants and Stimuli}

Ten participants were included in the analysis (7 healthy controls [HC], 3 females, age $22.7 \pm 2.5$ years; 3 color vision deficient [CVD], all male, including one deuteranope, one protanomalous, and one near-normal deutan serving as a specificity control).
Color vision was assessed using the 14-plate Ishihara test \cite{placeholder_ishihara1917}.
Given the small CVD sample, all CVD results are interpreted as individual case demonstrations, not population-level effects \cite{placeholder_crawford1998}.

Stimuli consisted of eight colors equally spaced from 0$^\circ$ to 315$^\circ$ in CIE~$L^*a^*b^*$ space, with lightness fixed at $L^* = 75$ and chroma distance of 40 from the achromatic point.
Participants performed a modified Rapid Serial Visual Presentation (RSVP) detection task \cite{placeholder_brouwer2009} across six runs per session, maintaining fixation while viewing 1.5~s color patches presented in optimized order.
Functional images were acquired on a Siemens 3T scanner (TR\,=\,1.5~s, 2~mm isotropic voxels, 24 oblique slices perpendicular to the calcarine sulcus), preprocessed with fMRIPrep, and projected to retinotopically-defined ROIs (V1, V2, hV4) using the Wang probabilistic atlas at 50\% threshold \cite{placeholder_wang2015}.
Response amplitudes were estimated using a two-stage FIR-then-GLM procedure \cite{placeholder_brouwer2009} and Procrustes-aligned across runs within each participant to improve cross-run consistency.

\subsection{Neural Phenotype: Interpolation Vulnerability}

We modeled hue-selective neural populations using half-wave rectified and squared sinusoidal basis functions \cite{placeholder_brouwer2009}: each of $K$ hypothetical channels was tuned to a preferred hue $\theta_k$, with response $r_k(\theta) = \max(0, \cos(\theta - \theta_k))^2$.
The weight matrix $W$ mapping channel outputs to voxel responses was estimated via ridge regression with generalized cross-validation (GCV) for regularization \cite{placeholder_golub1979}.

The target phenotype is the \emph{LOCO vulnerability profile}, a geometry-sensitive metric of representational structure.
Leave-one-color-out (LOCO) cross-validation tests whether the model can interpolate to novel hues not seen during training, assessing the preservation of continuous color-space geometry in the high-dimensional voxel pattern space: $W$ is trained on 7 of 8 colors and voxel responses to the held-out color are predicted.
The per-color prediction error defines an 8-element vulnerability vector $\mathbf{v} \in \mathbb{R}^8$ for each participant, capturing which hues are poorly interpolated from their neighbors---a structural deficit invisible to classification-based decoders.

This phenotype is specific to continuous geometry and distinct from color discrimination.
Leave-one-\emph{run}-out (LORO) cross-validation, which tests whether individual colors can be reliably distinguished across repetitions, shows no HC--CVD difference ($p = 0.668$), while LOCO interpolation error is significantly elevated in CVD.
Among the ROIs tested, the effect was most robust in hV4 ($p = 0.017$; $K{=}3$), consistent with hV4's established role in perceptual color representation \cite{placeholder_brouwer2009, placeholder_bannert2018}.
We therefore focus distortion fitting on hV4, with V1 analyses reported as cross-ROI validation in \cref{app:supporting}.

\subsection{Distortion Models}

We fit three forward models of hue distortion (\cref{fig:main}a), each parameterizing how a physical stimulus angle $\theta$ is transformed into a distorted angle $\theta'$:

\paragraph{Machado cone shift (1 DOF).}
A physiological model \cite{placeholder_machado2009} in which the anomalous cone's spectral peak is shifted by $\Delta\lambda \in [0, 20]$~nm.
This retinal-level model predicts hue distortion through altered L--M cone opponency.

\paragraph{Retinal + cortical gain (R+C; 2 DOF).}
Extends the cone shift with a cortical opponent-channel gain $g$:
\begin{equation}
rg' = rg_{\text{base}} + (1{+}g) \cdot (rg_{\text{ret}} - rg_{\text{base}})
\label{eq:rc}
\end{equation}
where $g{=}0$ reduces to pure Machado, $g{=}{-}1$ is exact cortical compensation, and $g{<}{-}1$ is overcompensation \cite{placeholder_tregillus2021}.

\paragraph{Fourier warp (4 DOF; overfitting ceiling).}
A model-free sinusoidal distortion field: $\delta(\theta) = a_1\!\sin\theta + b_1\!\cos\theta + a_2\!\sin 2\theta + b_2\!\cos 2\theta$.
With 4 parameters for 8 data points, this serves as a \emph{ceiling} to assess whether parsimonious models leave substantial variance unexplained.

\subsection{Fitting and Specificity}

For each model and candidate parameter(s), we simulate the LOCO vulnerability profile that a healthy observer would exhibit if viewing stimuli distorted by the model's forward map.
The fit is the Spearman rank correlation $\rho$ between simulated and observed vulnerability vectors, optimized over a grid ($\Delta\lambda$: 41 points; $g$: 25 points; Fourier: differential evolution).

Statistical significance is assessed by exact permutation of the 8-color labels ($8! = 40{,}320$ permutations).
A valid model must satisfy two criteria:
\begin{enumerate}
    \item \textbf{Sensitivity}: significant fit ($p < 0.05$) for CVD participants.
    \item \textbf{Specificity}: non-significant fit for the normal control (sub-10).
\end{enumerate}
A model that fits CVD subjects but also yields a false positive on the control is rejected as overfitting.

%==============================================================================
\section{Results}
\label{sec:results}
%==============================================================================

\subsection{Subject-Specific Model Selection}

\cref{tab:main} and \cref{fig:main} The central result is that distinct CVD participants require distinct low-dimensional distortion models, whereas the normal control yields no significant fitted distortion.

\begin{table}[t]
  \caption{Subject-specific distortion model fits (hV4 LOCO). Each CVD subject is best explained by a different parsimonious model. The normal control is correctly rejected by all models. Fourier serves as an overfitting ceiling.}
  \label{tab:main}
  \begin{center}
    \begin{small}
      \begin{tabular}{llcccc}
        \toprule
        Subject & Best model & DOF & $\rho$ & $p_{\text{perm}}$ \\
        \midrule
        Sub-08 (deutan) & R+C & 2 & 0.857 & \textbf{0.005**} \\
        Sub-09 (protan) & Machado & 1 & 0.762 & \textbf{0.018*} \\
        Sub-10 (control) & --- & --- & $-0.048$ & 0.559 \\
        \midrule
        \multicolumn{5}{l}{\emph{Overfitting ceiling (Fourier, 4 DOF):}} \\
        Sub-08 & Fourier & 4 & 0.976 & 0.0002 \\
        Sub-09 & Fourier & 4 & 0.762 & 0.018 \\
        \bottomrule
      \end{tabular}
    \end{small}
  \end{center}
  \vskip -0.15in
\end{table}

\textbf{Sub-09 (moderate protanomaly).}
The 1-DOF Machado model with $\Delta\lambda = 13.5$~nm---within the established range for moderate--severe protanomaly \cite{placeholder_machado2009}---captures the interpolation profile at $\rho = 0.762$, $p = 0.018$.
Adding a cortical gain parameter does not improve fit ($g = 0.0$ at optimum), indicating that a purely retinal model is sufficient for this participant.

\textbf{Sub-08 (mild deuteranomaly).}
The 1-DOF Machado model is only trend-level ($p = 0.058$).
The 2-DOF R+C model, with a small retinal shift ($\Delta\lambda = 2.0$~nm) and a cortical gain term ($g = +2.25$), achieves $\rho = 0.857$, $p = 0.005$.
The additional parameter is essential: the retinal shift alone explains little of the distortion in this participant's hV4 representation.

\textbf{Sub-10 (normal control).}
All parsimonious models yield $\Delta\lambda \approx 0$ and non-significant fits ($p = 0.559$), confirming specificity.

\textbf{Model parsimony matters.}
The 4-DOF Fourier model achieves $\rho = 0.976$ for sub-08 but is weakly constrained relative to the size of the observed profile (4 free parameters for 8 colors). For sub-09, its fit ($\rho = 0.762$) is identical to the 1-DOF Machado model, indicating that additional flexibility provides no benefit when the distortion is already captured by a parsimonious mechanistic model.

\begin{figure}[t]
  \vskip 0.2in
  \figplaceholder{\columnwidth}{%
    \textbf{Figure 1: Subject-Specific Distortion Profiles}\\[0.3em]
    \small \textbf{(a)} Model schematic: stimulus $\theta$ $\to$ retinal cone shift ($\Delta\lambda$) $\to$ optional cortical gain ($g$) $\to$ distorted $\theta'$ $\to$ forward encoding $\to$ LOCO evaluation.\\[0.3em]
    \textbf{(b--d)} 8-color LOCO vulnerability profiles.
    Gray bars: observed CVD profile.
    Colored line: best-model fit.
    \textbf{(b)} Sub-08: R+C fit ($\rho{=}0.857$, $p{=}0.005$). \textbf{(c)} Sub-09: Machado fit ($\rho{=}0.762$, $p{=}0.018$). \textbf{(d)} Sub-10: flat profile, $p{=}0.559$ (specificity pass).
  }
  \caption{The central result. Each CVD participant exhibits a distinct vulnerability profile (gray bars) in the high-dimensional voxel representation, captured by a different parsimonious model (colored lines). The normal control shows no structured vulnerability. Panel (a) illustrates the model hierarchy.}
  \label{fig:main}
  \vskip -0.1in
\end{figure}

\subsection{Fitted Parameters}

The fitted retinal parameters fall within established physiological ranges: sub-08's $\Delta\lambda = 2.0$~nm matches very mild deuteranomaly, and sub-09's $\Delta\lambda = 13.5$~nm matches moderate protanomaly \cite{placeholder_machado2009}.
Sub-09's cortical gain collapses to $g = 0$, consistent with a purely retinal account.

Sub-08's fitted cortical gain ($g = +2.25$) is less straightforward to interpret.
One possibility is that the mild retinal deficit ($\Delta\lambda = 2.0$~nm) produces a small spectral signal that cortical processing amplifies into a measurable interpolation deficit in hV4---an effect invisible under retinal-only modeling.
However, $g$ could alternatively be absorbing unmodeled variance (e.g., individual differences in voxel noise or attentional modulation) rather than reflecting a genuine cortical mechanism.
We treat this parameter as a provisional effective description of sub-08's distortion, not as a direct physiological measurement.
Cortical compensation magnitudes are known to vary widely across individuals \cite{placeholder_tregillus2021}, but the specific sign and magnitude observed here ($g > 0$, amplification rather than compensation) lacks direct precedent and requires independent validation.

\subsection{Geometric Collapse Determines Correction Limits}

Given a fitted forward model $D\!: \theta \to \theta'$, we ask whether an inverse correction exists---i.e., whether modified stimulus angles $\theta^*$ can be found such that $D(\theta^*) = \theta_{\text{target}}$ for all 8 colors.

\textbf{Sub-08 (mild, $\Delta\lambda = 2.0$~nm):} The R+C forward map preserves an opponent arc wide enough that all 8 hues remain separable.
Numerical pre-image search yields exact solutions ($< 0.001^\circ$ residual), producing an 8-point stimulus correction lookup table (\cref{fig:feasibility}a).

\textbf{Sub-09 (moderate, $\Delta\lambda = 13.5$~nm):} The Machado forward map compresses the 360$^\circ$ opponent arc to approximately 96$^\circ$, merging green, cyan, and blue into a single perceived angle ($\sim$282$^\circ$).
Exact inversion is structurally impossible for 4/8 colors (\cref{fig:feasibility}b).
A discriminability-oriented fallback (maximizing minimum pairwise angular separation) achieves 5.6$\times$ improvement over the unfiltered state but reaches only 48\% of the healthy-observer separation.

This severity-dependent boundary gives the forward model a practical role: the degree of geometric collapse in the fitted distortion field predicts whether exact stimulus-space correction is feasible for a given individual, before behavioral testing is conducted.

\begin{figure}[t]
  \vskip 0.2in
  \figplaceholder{\columnwidth}{%
    \textbf{Figure 2: Geometric Collapse and Correction Limits}\\[0.3em]
    \small \textbf{(a)} Sub-08 (mild): wide opponent arc; pre-image inversion restores all 8 colors ($<0.001^\circ$ error).\\
    \textbf{(b)} Sub-09 (moderate): arc compresses 360$^\circ$ $\to$ $\sim$96$^\circ$; three colors merge at $\sim$282$^\circ$.
    Exact inversion impossible.
    Separation-optimized fallback: $5.6\times$ improvement, 48\% of healthy separation.
  }
  \caption{The degree of geometric collapse in the fitted distortion field determines whether exact stimulus-space correction is achievable. Mild distortion preserves separability and admits exact inversion; moderate distortion collapses multiple hues onto one perceived angle, precluding exact recovery.}
  \label{fig:feasibility}
  \vskip -0.1in
\end{figure}

%==============================================================================
\section{Discussion}
\label{sec:discussion}
%==============================================================================

\paragraph{Inferring distortion from complex signals.}
The core methodological contribution is that subject-specific, low-dimensional distortion fields (1--2 parameters) can be inferred from high-dimensional fMRI voxel patterns (hundreds of voxels $\times$ 8 colors $\times$ 6 runs).
The LOCO vulnerability profile serves as the bridge: it compresses the high-dimensional representation into a geometry-sensitive 8-element summary that is sensitive to continuous structure rather than categorical boundaries.
This approach may generalize beyond CVD to other conditions where sensory representations are distorted rather than absent.

\paragraph{What this framework does and does not establish.}
We have shown that subject-specific interpolation profiles can be recovered by parsimonious forward models, with correct specificity rejection and physiologically plausible retinal parameters.
We have \emph{not} shown that the resulting corrections improve perceptual discrimination---that requires behavioral validation, which is future work.
The current evidence supports the claim that the method identifies individual distortion profiles and estimates whether stimulus-space correction is feasible under the fitted model; it does not yet establish that these corrections improve perception.

\paragraph{Why different models for different subjects?}
Sub-09's profile is captured by pure cone shift ($g = 0$) while sub-08 requires the additional cortical gain term.
This heterogeneity is itself an important result: it suggests that the neural consequences of CVD cannot be reduced to diagnostic category alone, but instead reflect person-specific variation in retinal severity and downstream processing that is only visible under individualized modeling.

\paragraph{Why reject more flexible models?}
The Fourier model ($\rho = 0.976$ for sub-08) fits nearly perfectly but affords poor generalization guarantees with 4 parameters for 8 data points.
The critical test is not fit quality but the joint requirement of sensitivity \emph{and} specificity: a valid model must explain CVD profiles while rejecting the normal control.
Under this criterion, parsimonious models are the meaningful result, because they constrain the solution space enough that a null participant cannot produce a spurious fit.

\paragraph{Limitations.}
Three CVD participants are insufficient for population-level claims; this work is a proof-of-concept for individualized distortion modeling.
The fMRI requirement limits near-term scalability.
The hV4 focus, while motivated by the strongest geometry-sensitive HC–CVD separation under LOCO, means that distortion structure in earlier visual areas (V1, V2) is not directly modeled in the main analysis (see \cref{app:supporting} for cross-ROI evidence).
The interpolation phenotype depends on forward encoding model specification (basis type, number of channels, regularization), though prior analyses confirm robustness across reasonable parameter choices.

\paragraph{Relation to prior work.}
Our approach extends Brouwer \& Heeger's \cite{placeholder_brouwer2009} forward encoding framework---which demonstrated LOCO reconstruction in hV4 for healthy observers---to CVD with individual-level distortion fitting.
The distinction between preserved classification and impaired interpolation is consistent with reports that CVD individuals maintain categorical but not continuous color representations \cite{placeholder_bannert2018}.
The cortical compensation parameters relate to accumulating evidence for long-term neural adaptation in anomalous trichromacy \cite{placeholder_tregillus2021, placeholder_emery2021}.

%==============================================================================
\section{Conclusion}
\label{sec:conclusion}
%==============================================================================

We present a proof-of-concept framework for inferring subject-specific distortion fields from high-dimensional neural representations in color vision deficiency.
The framework extracts a geometry-sensitive interpolation phenotype from fMRI data, fits it with parsimonious mechanistic models, and predicts whether stimulus-space correction is feasible for a given individual based on the degree of geometric collapse.
The critical next step is behavioral validation: determining whether the model-prescribed corrections actually improve color discrimination in the fitted individuals.

%==============================================================================
\section*{Impact Statement}
%==============================================================================

This work develops personalized neural models of color vision deficiency with the goal of informing assistive technology design.
Potential positive impacts include individually-tailored color corrections for digital displays and accessibility tools.
The framework currently requires fMRI data, limiting immediate clinical deployment.
Behavioral validation and informed consent are essential prerequisites before any correction is applied, as a misspecified model could impair rather than improve visual function.

%==============================================================================
\bibliography{genbio_draft}
\bibliographystyle{icml2026}

%==============================================================================
%==============================================================================
\newpage
\appendix
\onecolumn

\section{Supporting Analyses}
\label{app:supporting}

\subsection{Cross-ROI Validation: V1 Geometric Evidence}

Beyond the primary hV4 LOCO fitting (main text), we performed an independent geometric validation using pairwise representational distance matrices (RDMs) in V1.
A 2-component angular dilation model---parameterized by S-cone expansion $\beta_s$ and confusion-axis modulation $\beta_c$---was fitted to the voxel-space $\Delta\text{RDM} = \text{RDM}_{\text{CVD}} - \overline{\text{RDM}}_{\text{HC}}$.
Sub-09 achieved crossnobis cosine $p = 0.007$ with $\beta_s = 23.0^\circ \pm 10.2^\circ$ (bootstrap CI excluding 0).
The cross-subject mean $\beta_s \approx 21.5^\circ$ matches the $21.4^\circ$ B-Y axis rotation independently reported from behavioral hue-scaling \cite{placeholder_emery2021}---convergent evidence from an entirely different fitting criterion and ROI.

This supplementary analysis confirms that the distortion captured by hV4 LOCO is not ROI-specific: it is also visible as geometric distortion in V1 pairwise distance structure, though the optimal model and parameter regime differ across ROIs.

\subsection{Model Comparison: Full Table}

\begin{table}[h]
  \caption{Full model comparison across hV4 LOCO and V1 $\Delta$RDM criteria.}
  \label{tab:full}
  \begin{center}
    \begin{small}
      \begin{tabular}{llccccc}
        \toprule
        Subject & Model & DOF & hV4 LOCO $\rho$ & hV4 LOCO $p$ & V1 $\Delta$RDM & Notes \\
        \midrule
        \multirow{4}{*}{Sub-08} & Machado & 1 & 0.619 & 0.058 & anti-corr. & Trending only \\
        & R+C & 2 & \textbf{0.857} & \textbf{0.005**} & 0.179 & Primary \\
        & 2-Comp & 2 & 0.881 & 0.004** & CI excl.\ 0 & Alternative \\
        & Fourier & 4 & 0.976 & 0.0002 & --- & Ceiling \\
        \midrule
        \multirow{4}{*}{Sub-09} & Machado & 1 & \textbf{0.762} & \textbf{0.018*} & $+0.091$ & Primary \\
        & R+C & 2 & 0.762 & 0.018 & 0.026* & $g{=}0$ (= Machado) \\
        & 2-Comp & 2 & 0.690 & 0.035* & 0.007*** & Cross-criterion \\
        & Fourier & 4 & 0.762 & 0.018 & --- & No gain \\
        \midrule
        Sub-10 & All & --- & $-0.048$ & 0.559 & --- & Specificity pass \\
        \bottomrule
      \end{tabular}
    \end{small}
  \end{center}
\end{table}

\subsection{Correction Feasibility Details}

\paragraph{Sub-08 (R+C pre-image).}
All 8 colors admit exact pre-image solutions (mean $|\delta| = 23.5^\circ$, max $= 42.9^\circ$, residual $< 0.001^\circ$).
The resulting 8-point lookup table is directly implementable on any digital display.

\paragraph{Sub-09 (Machado pre-image).}
At $\Delta\lambda = 13.5$~nm, the opponent-space arc compresses from 360$^\circ$ to $\sim$96$^\circ$.
Colors c4 (green), c5 (cyan), and c6 (blue) converge to a single perceived angle ($\sim$282$^\circ$).
Exact pre-image: 4/8 colors.
Separation-optimized fallback: min separation $1.03^\circ \to 5.76^\circ$ ($5.6\times$), mean separation $39.3^\circ \to 24.3^\circ$ (healthy: 70.7$^\circ$).

\end{document}
